% Preface: the README's About the Course and Written Examinations sections, the course plan and
% lecture content from info/README.md, and what the lectures take for granted, as the front matter
% of a book.

\chapter{Preface}
\markboth{Preface}{Preface}

This book covers the fundamentals of machine learning, and it covers them by building them. Every
model in it is implemented from scratch in C++, without an external machine learning library, so
that nothing happens inside it that you have not written yourself:
\begin{itemize}
  \item Linear regression and gradient-based optimization.
  \item Neural networks: feedforward, backpropagation, and activation functions.
  \item Convolutional neural networks (CNNs) for image classification: kernels, pooling, and
    flatten layers.
  \item Implementing machine learning algorithms from scratch in C++.
\end{itemize}

\textbf{Along the way you will train models by hand as well as in software.} Every algorithm the
book implements is first worked through on paper, with numbers small enough to follow, and then
written as code that reproduces those numbers. The code you write adds up to:
\begin{itemize}
  \item A linear regression model trained with a fixed learning rate, and a second one that picks
    its own.
  \item A simple neural network, built on a dense layer written from scratch.
  \item Convolutional, pooling, and flatten layers, first as simple structs and then as classes
    behind an interface, wired together into a network that tells four hand-drawn digits apart.
\end{itemize}

\section*{What you should be able to do at the end}
Having worked through the book, you should be able to:
\begin{itemize}
  \item Implement basic machine learning algorithms from scratch in C++.
  \item Train and evaluate neural networks, both by hand and in software.
  \item Reason about design choices and limitations.
\end{itemize}

\section*{What you are assumed to know}
The subject of this book is \textbf{machine learning}, not C++. It introduces
\code{std::vector}, the one container every model in it is built on, in \secref{c1:sec:vectors},
and uses the rest of C++17 without introducing it: classes and structs, public inheritance and
pure virtual interfaces, references and \code{const}, the keywords \code{explicit},
\code{noexcept}, \code{final}, \code{override} and \code{[[nodiscard]]}, deleted copy and move
operations, and, from Chapter~\ref{c8:ch:conv} on, \code{std::unique_ptr} and a factory. You should
also be at home in a terminal, compiling a program with \code{g++} and building one with
\code{make}.

The mathematics asks for less than it may look. Everything the book computes is sums, products and
differences, plus the derivatives of the three activation functions it implements, and every one of
those derivatives is given. A calculator is all the hand calculations need.

\section*{How the book is organised}
The book is typeset from a ten-lecture course, and each chapter is one lecture. The course is taught
in two parts with a natural break between them, and the book keeps the break: Part~I builds up from
linear regression to a neural network with a dense layer written from scratch, and Part~II builds a
convolutional neural network one layer type at a time, closing with a review that ties linear
regression, neural networks, and CNNs together.

\begin{booktable}{@{}p{5.2em}p{3.6em}L@{}}
  \thead{Chapter} & \thead{Lecture} & \thead{Topic} \\
  \midrule
  \multicolumn{3}{@{}l}{\textit{Part I: Linear Regression and Neural Networks}} \\
  1 & \file{L01} & Linear regression (part I) \\
  2 & \file{L02} & Linear regression (part II) \\
  3 & \file{L03} & Neural networks (part I) \\
  4 & \file{L04} & Neural networks (part II) \\
  5 & \file{L05} & Dense layer \\
  \midrule
  \multicolumn{3}{@{}l}{\textit{Part II: Convolutional Neural Networks}} \\
  6 & \file{L06} & Convolutional neural networks (part I) \\
  7 & \file{L07} & Convolutional neural networks (part II) \\
  8 & \file{L08} & Convolutional neural networks (part III) \\
  9 & \file{L09} & Convolutional neural networks (part IV) \\
  10 & \file{L10} & Flatten layer and course review \\
\end{booktable}

A chapter's sections are its lecture's appendices, in the same order and with the same
subsections, so a chapter here and a lecture directory in the repository can be read side by side.
Every chapter opens with what is in it, works through the material, closes with a review of what
you should now be able to explain, and ends with its exercises. The appendices hold the worked
solutions to the exercises answered on paper, and the two written papers described below.

\section*{What the exercises look like}
The exercises come in \textbf{sets}. A set builds one program, or one part of a larger one, or
works through one calculation, and each exercise in a set adds a piece of it. Every exercise is
labelled with its kind:

\begin{booktable}{@{}p{7.4em}L@{}}
  \thead{Kind} & \thead{What it asks of you} \\
  \midrule
  \kindtag{Code} & Write C++, compile it, and check it against the lecture's test suite, which
    specifies the behaviour the exercise asks for, case by case. \\
  \kindtag[fillaccent3]{Calculation} & Work it out on paper: train a model by hand, trace an image
    through a layer, or count what a network has to store. \\
  \kindtag[fillaccent2]{Reflection} & Answer in prose: recognise a layer from its algorithm, or
    decide whether a claim about one is true, and say why. \\
\end{booktable}

\textbf{The solutions to the Calculation and Reflection exercises are in this book}, in
Appendix~\ref{sol:ch:solutions}, worked through step by step, and every such exercise ends with the
page its solution is on. Work each one before you turn to it; a hand calculation you have only read
teaches you much less than one you got wrong first.

\textbf{The solutions to the Code exercises are not in this book.} They are published in the
companion repository, lecture by lecture, and each lecture's README says where. Clone the repository
(see \emph{Setting up} below) and you can run a reference program beside your own.

\section*{Two written papers, and what they are not}
Nothing in the course this book comes from is marked. Assessment is the exercises after every
chapter, the test suite most of them ship with, and the self-assessment exercises in
\secref{c10:sec:exercises}, each of which checks a component you built or a number you computed
beside it.

Appendices~\ref{exam:ch:about} to \ref{ansb:ch:answers} hold two four-hour papers with worked
solutions, and they check something else: \textbf{your theoretical knowledge, on paper, with
nothing in front of you.} Nine questions each, mixing theory with code, and every derivation the
chapters make asked for without the chapter open beside you.

\textbf{They exist purely so that you can test your own knowledge after working through the book.
They gate nothing, they are not a qualification, and no part of the book requires them.}

\textbf{Take one after the last chapter.} Both papers draw on all ten chapters, so sitting one
partway through examines material you have not read yet.

\section*{Setting up}
Everything in this book runs in a Linux terminal: on Linux itself, or on Windows through WSL. You
need three things:
\begin{itemize}
  \item \textbf{GCC and G++}, with C++17 support, to compile the programs.
  \item \textbf{Linux or WSL}, to compile and run them in a terminal environment.
  \item \textbf{An editor}; Visual Studio Code is the one the course uses, available at
    \url{https://code.visualstudio.com/download}.
\end{itemize}
Chapter~\ref{c7:ch:layers} also points to reference implementations of its layers in Python, which
run with \code{python3}. For a small program you want to try without installing anything, an online
compiler such as \url{https://www.onlinegdb.com/online_c_compiler} works too, with the language
switched to C++17.

Then clone the companion repository, together with the test framework it keeps as a submodule:

\begin{shell}
git clone --recursive https://github.com/qrtech-academy/machine-learning.git
\end{shell}

\noindent (In a clone made without \code{--recursive}, run \code{git submodule update --init}
once.) Three commands at its root cover the repository as a whole:

\begin{shell}
make build               # Build every lecture demo that has a Makefile.
make test                # Build and run every test suite that has a solution to run against.
make format-check        # Check the formatting of every C++ and Python file.
\end{shell}

\noindent You write your exercise solutions in the repository, and each chapter's exercises say
where. Each lecture's test suite already knows where that is and runs with a single \code{make}:
its tests use nothing but the public interface the exercises specify, so any implementation that
follows the specification passes, however it is written on the inside, and a test that fails names
the assertion, the two values involved, and the file and line it came from.
